%======================================================================
% Section 13: Applications of Multiphysics Wave Mechanics
%======================================================================
\section{Applications of Multiphysics Wave Mechanics}
\label{sec:applications}

This section reviews applications of multiphysics wave mechanics across engineering and science. Each application is assessed for peer-reviewed support, providing an evidence map of which applications have strong theoretical and experimental foundations.

\subsection{Smart structures and adaptive materials}
\label{sec:13.1}

Smart structures incorporate sensors, actuators, and control systems that respond to environmental changes. Piezothermoelastic materials are used in:
\begin{itemize}
\item Adaptive aerospace structures (wing shaping, vibration control),
\item Active vibration suppression in precision machinery,
\item Structural health monitoring with embedded piezoelectric sensors.
\end{itemize}

\subsubsection{Evidence}
Peer-reviewed studies demonstrate piezothermoelastic coupling in smart structures, with experimental validation of sensor-actuator performance.

\subsection{Nondestructive evaluation and sensing}
\label{sec:13.2}

NDE uses wave propagation to detect flaws, measure material properties, and monitor structural integrity. Applications include:
\begin{itemize}
\item Ultrasonic testing for crack detection,
\item Thermal wave imaging for subsurface defect detection,
\item Acoustic emission monitoring for damage assessment.
\end{itemize}

\subsubsection{Evidence}
Extensive peer-reviewed literature supports ultrasonic NDE; thermal wave imaging has growing experimental support.

\subsection{Energy harvesting and conversion}
\label{sec:13.3}

Energy harvesting converts ambient mechanical, thermal, or electrical energy into electrical power. Piezothermoelectric energy harvesters use:
\begin{itemize}
\item Vibration-induced piezoelectric power generation,
\item Thermal gradient-driven thermoelectric conversion,
\item Coupled thermo-piezoelectric conversion.
\end{itemize}

\subsubsection{Evidence}
Peer-reviewed studies demonstrate piezoelectric energy harvesting; coupled thermo-piezoelectric conversion has emerging experimental support.

\subsection{Geophysics and seismology}
\label{sec:13.4}

Geophysical applications include:
\begin{itemize}
\item Seismic wave propagation in the Earth's crust,
\item Thermal wave propagation in geothermal systems,
\item Poroelastic wave propagation in saturated reservoirs.
\end{itemize}

\subsubsection{Evidence}
Seismic wave theory is well established; poroelastic wave propagation has strong theoretical and experimental support.

\subsection{Biomedical applications (ultrasound, thermal therapy)}
\label{sec:13.5}

Biomedical applications include:
\begin{itemize}
\item Medical ultrasound imaging and therapy,
\item Thermal ablation for cancer treatment,
\item Photoacoustic imaging,
\item Temperature-sensitive drug delivery.
\end{itemize}

\subsubsection{Evidence}
Medical ultrasound is well established; thermal ablation has growing clinical evidence; photoacoustic imaging is an active research area.

%======================================================================
% Table T8: Applications with evidence levels
%======================================================================
\begin{table}[htbp]
\centering
\small
\caption{Evidence map of applications across theory families: evidence level and representative references.}
\label{tab:T8_applications}
\begin{tabularx}{\textwidth}{@{}lXXXX@{}}
\toprule
\textbf{Application} & \textbf{Theory families} & \textbf{Evidence level} & \textbf{Representative references} \\
\midrule
Smart structures / adaptive aerospace & Piezothermoelastic & High & Multiple experimental studies on sensor--actuator systems \\
Nondestructive evaluation (ultrasonic) & Classical, LS, GL & High & Extensive clinical and industrial literature \\
Nondestructive evaluation (thermal wave) & LS, DPL, GN & Medium & Growing experimental literature on thermal wave imaging \\
Energy harvesting (piezoelectric) & Piezoelectric & High & Multiple experimental studies on vibration energy harvesting \\
Energy harvesting (thermo-piezoelectric) & Piezothermoelastic & Medium & Emerging experimental support \\
Geophysics / seismology & Classical, poroelastic & High & Well-established seismic wave theory \\
Poroelastic reservoir characterization & Double porous & Medium & Theoretical + field studies on hydrocarbon reservoirs \\
Biomedical ultrasound imaging & Classical & High & Clinical practice; extensive validation \\
Thermal ablation (cancer treatment) & Classical, LS & Medium & Growing clinical evidence; treatment planning \\
Photoacoustic imaging & LS, DPL & Medium & Active research area; preclinical studies \\
Nanoscale devices / MEMS & Nonlocal & Medium & Theoretical + some experimental validation \\
Fractional viscoelastic materials & Fractional, MDD & Medium & Polymer and biological tissue studies \\
\bottomrule
\end{tabularx}
\end{table}


\begin{figure}[htbp]
\centering
\includegraphics[width=0.95\textwidth]{fig12-application_evidence.pdf}
\caption{Application evidence map: peer-reviewed evidence levels across applications and theory families (Section 13.6).}
\label{fig:evidence_map}
\end{figure}

\subsection{Evidence map: which applications have peer-reviewed support}
\label{sec:13.6}

\begin{table}[h]
\centering
\small
\caption{Evidence map of applications across theory families.}
\label{tab:evidence_map}
\begin{tabularx}{\textwidth}{@{}lXrX@{}}
\toprule
\textbf{Application} & \textbf{Theory families} & \textbf{Evidence level} & \textbf{Key references} \\
\midrule
Smart structures & Piezothermoelastic & High & Multiple experimental studies \\
NDE (ultrasonic) & Classical, LS, GL & High & Extensive literature \\
NDE (thermal wave) & LS, DPL, GN & Medium & Growing literature \\
Energy harvesting & Piezoelectric & High & Multiple experimental studies \\
Geophysics/seismology & Classical, poroelastic & High & Well-established theory \\
Biomedical ultrasound & Classical & High & Clinical practice \\
Thermal ablation & Classical, LS & Medium & Growing clinical evidence \\
Photoacoustic imaging & LS, DPL & Medium & Active research area \\
Nanoscale devices & Nonlocal & Medium & Theoretical + some experimental \\
\bottomrule
\end{tabularx}
\end{table}

\subsection{Notation summary for Section 13}

\begin{table}[h]
\centering
\small
\caption{Notation introduced in Section 13.}
\label{tab:notation-sec13}
\begin{tabularx}{\textwidth}{@{}lXlr@{}}
\toprule
\textbf{Symbol} & \textbf{Meaning} & \textbf{Units} & \textbf{Equation} \\
\midrule
--- & --- & --- & --- \\
\bottomrule
\end{tabularx}
\end{table}

\subsection{Open questions for Section 13}

\begin{enumerate}[label=\textbf{OQ13.\arabic*}]
\item Which applications have the strongest experimental validation across theory families? \label{oq:13.1}
\item What are the barriers to translating generalized thermoelasticity theories into practical applications? \label{oq:13.2}
\end{enumerate}

% LocalWords: 
